% !TeX root = ../../main.tex
\section{两体问题}

\begin{solution}[变换到质心系]
\problem{利用变量变换从 \((2.3.5)\) 式推导 \((2.3.6)\) 式。}
\answer
原来的两体问题拉格朗日量为
\[
\mathcal{L}
=\frac{1}{2}m_1|\dot{\mathbf r}_1|^2+\frac{1}{2}m_2|\dot{\mathbf r}_2|^2
-V(\mathbf r_1-\mathbf r_2).
\]
为了将其改写成质心运动和相对运动的形式，引入质心坐标 \(\mathbf r_{\rm cm}\) 与相对坐标 \(\mathbf r\)：
\[
\mathbf r_{\rm cm}=\frac{m_1\mathbf r_1+m_2\mathbf r_2}{m_1+m_2},\qquad
\mathbf r=\mathbf r_1-\mathbf r_2.
\]
由这两个定义反解出 \(\mathbf r_1\) 和 \(\mathbf r_2\)。设 \(M=m_1+m_2\)，则
\[
\mathbf r_1=\mathbf r_{\rm cm}+\frac{m_2}{M}\mathbf r,\qquad
\mathbf r_2=\mathbf r_{\rm cm}-\frac{m_1}{M}\mathbf r.
\]
对时间求导，得到
\[
\dot{\mathbf r}_1=\dot{\mathbf r}_{\rm cm}+\frac{m_2}{M}\dot{\mathbf r},\qquad
\dot{\mathbf r}_2=\dot{\mathbf r}_{\rm cm}-\frac{m_1}{M}\dot{\mathbf r}.
\]
于是两个速度平方分别为
\[
|\dot{\mathbf r}_1|^2
=
|\dot{\mathbf r}_{\rm cm}|^2
+\frac{2m_2}{M}\dot{\mathbf r}_{\rm cm}\cdot\dot{\mathbf r}
+\frac{m_2^2}{M^2}|\dot{\mathbf r}|^2,
\]
\[
|\dot{\mathbf r}_2|^2
=
|\dot{\mathbf r}_{\rm cm}|^2
-\frac{2m_1}{M}\dot{\mathbf r}_{\rm cm}\cdot\dot{\mathbf r}
+\frac{m_1^2}{M^2}|\dot{\mathbf r}|^2.
\]
代入原拉格朗日量，得
\[
\begin{aligned}
\mathcal{L}
={}&\frac{1}{2}m_1\left(
|\dot{\mathbf r}_{\rm cm}|^2
+\frac{2m_2}{M}\dot{\mathbf r}_{\rm cm}\cdot\dot{\mathbf r}
+\frac{m_2^2}{M^2}|\dot{\mathbf r}|^2
\right)\\
&+\frac{1}{2}m_2\left(
|\dot{\mathbf r}_{\rm cm}|^2
-\frac{2m_1}{M}\dot{\mathbf r}_{\rm cm}\cdot\dot{\mathbf r}
+\frac{m_1^2}{M^2}|\dot{\mathbf r}|^2
\right)-V(\mathbf r).
\end{aligned}
\]
整理各项。首先，质心速度平方项为
\[
\frac{1}{2}m_1|\dot{\mathbf r}_{\rm cm}|^2+\frac{1}{2}m_2|\dot{\mathbf r}_{\rm cm}|^2
=\frac{1}{2}(m_1+m_2)|\dot{\mathbf r}_{\rm cm}|^2.
\]
其次，交叉项为
\[
\frac{1}{2}m_1\frac{2m_2}{M}\dot{\mathbf r}_{\rm cm}\cdot\dot{\mathbf r}
+\frac{1}{2}m_2\left(-\frac{2m_1}{M}\dot{\mathbf r}_{\rm cm}\cdot\dot{\mathbf r}\right)
=0,
\]
因此质心运动和相对运动没有速度交叉项。最后，相对速度平方项为
\[
\frac{1}{2}m_1\frac{m_2^2}{M^2}|\dot{\mathbf r}|^2
+\frac{1}{2}m_2\frac{m_1^2}{M^2}|\dot{\mathbf r}|^2
=\frac{1}{2}\frac{m_1m_2(m_1+m_2)}{M^2}|\dot{\mathbf r}|^2
=\frac{1}{2}\frac{m_1m_2}{m_1+m_2}|\dot{\mathbf r}|^2.
\]
因此拉格朗日量化为
\[
\mathcal{L}
=
\frac{1}{2}(m_1+m_2)|\dot{\mathbf r}_{\rm cm}|^2
+\frac{1}{2}\frac{m_1m_2}{m_1+m_2}|\dot{\mathbf r}|^2
-V(\mathbf r).
\]
若定义总质量 \(M=m_1+m_2\) 和约化质量
\[
\mu=\frac{m_1m_2}{m_1+m_2},
\]
则上式可写为更紧凑的形式
\[
\mathcal{L}
=
\frac{1}{2}M|\dot{\mathbf r}_{\rm cm}|^2
+\frac{1}{2}\mu|\dot{\mathbf r}|^2
-V(\mathbf r).
\]
这就是由 \((2.3.5)\) 经过变量变换得到的 \((2.3.6)\) 式，表明两体问题可以分解为一个质量为 \(M\) 的自由质心运动，以及一个质量为 \(\mu\)、处在势 \(V(\mathbf r)\) 中的相对运动。
\end{solution}

